\begin{table}[!ht]
  \centering
  \footnotesize
  \setlength{\tabcolsep}{4pt}
  \begin{threeparttable}
    \caption{O que cada propriedade do esquema do kw faz com a exportação e o
      que foi feito.}
    \label{tab:esquema}
    \begin{tabular}{@{}>{\raggedright\arraybackslash}p{4.0cm}
                      >{\raggedright\arraybackslash}p{3.8cm}
                      >{\raggedright\arraybackslash}p{6.4cm}@{}}
      \toprule
      Propriedade do esquema & Efeito na exportação & O que foi feito \\
      \midrule
      identificador por ordem de inserção, sem garantia de não reaproveitamento
        & um evento novo poderia passar por um já enviado
        & contornado: a marca confere identificador e conteúdo, e a chave de
          repetição usa o identificador. Proposta: impedir o reaproveitamento,
          o que exige recriar a tabela \\
      \addlinespace
      data e hora em texto, na hora local, sem fuso
        & a virada da semana fica incerta
        & contornado pela granularidade de semana. Proposta: uma coluna com o
          instante em UTC \\
      \addlinespace
      sem versão do esquema
        & uma coluna nova depende de verificação no ponto de uso
        & feito o mínimo: uma verificação antes do envio, que só acrescenta o que
          falta. Proposta: versão do esquema e uma guarda de migração \\
      \addlinespace
      sem estado de exportação
        & reenvio e contagem dupla
        & feito: uma coluna no evento, com um índice, marca o que foi enviado \\
      \bottomrule
    \end{tabular}
    \begin{tablenotes}
      \footnotesize
      \item Justificativas na Seção~\ref{sec:banco-local}.
    \end{tablenotes}
  \end{threeparttable}
\end{table}
